\documentclass[10pt,a4paper]{amsart}
\usepackage[margin=19mm]{geometry}
\usepackage{amsmath,amssymb,mathtools}
\usepackage[scr=rsfs,cal=euler]{mathalfa}
\usepackage{xfrac}
\usepackage[unicode,hidelinks,bookmarks=false]{hyperref}
\newcommand{\ie}{i.e.\ }
\newcommand{\cf}{cf.\ }
\newcommand{\resp}{respectively\ }
\newcommand{\Subsection}{\subsection}
\providecommand{\nobreakdash}{\nobreak\hbox{-}}
\setlength{\emergencystretch}{2em}
\multlinegap0pt
\usepackage{dsfont}
\usepackage{amssymb}
\usepackage{enumerate}
\newtheorem{hyp}{Hypothesis}
\newtheorem{lem}{Lemma}
\newcommand{\R}{\mathbb{R}}
\title{\vspace{-.4in}{\fontsize{18}{18}\selectfont \textbf{Saddle Transport and Chaos in the Double Pendulum}}\vspace{-.15in}}
\author{Kadierdan Kaheman, Jason J. Bramburger, J. Nathan Kutz, Steven L. Brunton}
\date{Source: arXiv:2209.10132v2 (CC BY 4.0)}
\begin{document}
\maketitle

\noindent\emph{Source and licence.} \vspace{-.4in}{\fontsize{18}{18}\selectfont \textbf{Saddle Transport and Chaos in the Double Pendulum}}\vspace{-.15in}. Version arXiv:2209.10132v2 (CC BY 4.0), distributed by arXiv under the Creative Commons Attribution 4.0 International licence, http://creativecommons.org/licenses/by/4.0/, as recorded in the arXiv licence metadata for this exact version and shown on the abstract page https://arxiv.org/abs/2209.10132v2. Full licence notice: \texttt{licenses/arxiv-2209.10132-CC-BY-4.0.txt}.\par\smallskip

\noindent\textit{Editorial scope.} Counted unit: the article's lem lem:Shil (source0.tex lines 826--835). The article's complete original proof of it is delivered with it (source0.tex lines 837--850, one \texttt{proof} environment of 14 lines). No mathematics is written, repaired or re-derived: the statement and the proof are the article's own lines, in the article's own notation, and no retained line is substituted or re-flowed.  Retained uncounted as the source-local context the proof needs: lines 657--659; lines 662--664; lines 668--674; lines 821--823.  Nothing was omitted from any retained span, so the package carries no omission record.  No retained line contains a citation, so the wrapper carries no bibliography and no bibliography entry.  No retained line contains a figure environment, an image-inclusion command or drawing code, so no image is delivered and the package declares no asset dependency.  Editorial formatting only, in addition: an AMS \texttt{amsart} preamble that restates the article's own declarations and macros used by the retained lines, namely the environments \texttt{hyp}, \texttt{lem} on separate counters, together with the standard packages those lines need.  The retained environments print in this document as Hyp 1, Lemma 1 in order of appearance, whereas the article prints the same items with the numbering of its own full document.  The complete native source of the article as distributed in the arXiv e-print for this exact version is bundled unchanged as \texttt{sources/130999/source0.tex}.
\begin{equation}\label{ODE}
	\dot{u} = F(u),
\end{equation}

\begin{hyp}\label{hyp:Ham} %Hypothesis: Hamiltonian structure
There exists a smooth function $\mathcal{H}:\R^4\to\R$ such that $\langle F(u),\nabla \mathcal{H}(u) \rangle = 0$ for all $u \in \R^4$.
\end{hyp}

\begin{hyp}\label{hyp:Periodic} %Hypothesis: Periodic orbits
There exists $p \geq 1$ such that~\eqref{ODE} has periodic orbits $\gamma_1(t), \gamma_2(t), \dots, \gamma_p(t)$ satisfying 
\begin{equation}
	\mathcal{H}(\gamma_i(t)) = \mathcal{H}(\gamma_j(t))
\end{equation}
for all $1 \leq i,j \leq p$ and all $t \in \R$. Moreover, each $\gamma_j$ is hyperbolic and so has precisely two Floquet multipliers at one and no others on the unit circle.
\end{hyp}

\begin{equation}\label{Psec}
	P_i: \Sigma'_i \to \Sigma_i,
\end{equation} 

\begin{lem}\label{lem:Shil} %Lemma: Shilnikov variables
Assume Hypotheses~\ref{hyp:Ham} and \ref{hyp:Periodic} are met. Then there exists $\delta > 0$ such that for each ${i \in \{1,\dots,k+1\}}$ there is a smooth change of coordinates mapping $u\in\Sigma'_i$ to $v_i = (v^s_i,v^u_i)$ near the fixed point $u = \tilde\gamma_i(0)$, and smooth functions $f^s_{i,j},f^u_{i,j}:\mathcal{I}\times\mathcal{I}\to \mathbb{R}$, $j = 1,2$, so that (\ref{Psec}) is of the form 
	\begin{equation}\label{Shil}
		\begin{split}
			v^s_{i,n+1} &= [\lambda_i^{-1} + f_{i,1}^s(v^s_{i,n},v^u_{i,n})v^s_{i,n} + f_{i,2}^s(v^s_{i,n},v^u_{i,n})v^u_{i,n}]v^s_{i,n}, \\
			v^u_{i,n+1} &= [\lambda_i + f_{i,1}^u(v^s_{i,n},v^u_{i,n})v^s_{i,n} + f_{i,2}^u(v^s_{i,n},v^u_{i,n})v^u_{i,n}]v^u_{i,n}, \\
		\end{split}
	\end{equation}	
	where $v^s_{i,n},v^u_{i,n} \in \mathcal{I} := [-\delta,\delta]$.
\end{lem}

\begin{proof}
Let $\xi^s_i$ be the eigenvector of $DP_i(\tilde\gamma_i(0))$ associated to $\lambda_i^{-1}$ and $\xi^u_i$ be the eigenvector associated to $\lambda_i$. It follows from the stable manifold theorem for maps that there exists a $\delta_i > 0$ and smooth functions $w^s_i,w^u_i:[-\delta_i,\delta_i] \to \mathbb{R}$ with $w^s_i(0) = (w^s_i)'(0) = w^u_i(0) = (w^u_i)'(0) = 0$ such that the local stable and unstable manifolds of $\tilde\gamma_i(0)$ can be written 
\begin{equation}
	\begin{split}
		W^s_\mathrm{loc}(\tilde\gamma_i(0)) &= \{\tilde\gamma_i(0) + v^s_i\xi^s_i + w^s_i(v^s_i)\xi^u_i:\ v^s_i \in [-\delta_i,\delta_i]\}, \\
		W^u_\mathrm{loc}(\tilde\gamma_i(0)) &= \{\tilde\gamma_i(0) + v^u_i\xi^u_i + w^u_i(v^u_i)\xi^s_i:\ v^u_i \in [-\delta_i,\delta_i]\}.
	\end{split}
\end{equation}
Then, for $u\in\Sigma_i'$ in a small neighborhood of $\tilde\gamma_i(0)$ we may introduce the change of variable
\begin{equation}
	(v^s_i,v^u_i) \mapsto u = \tilde\gamma_i(0) + (v^s_i\xi^s_i + w^s_i(v^s_i)\xi^u_i) + (v^u_i\xi^u_i + w^u_i(v^u_i)\xi^s_i). 
\end{equation}
From the tangency properties of the $w^s_i,w^u_i$ functions at $(v^s_i,v^u_i) = (0,0)$, it follows that the above change of variable is a local diffeomorphism. Putting this change of variable into the Poincar\'e map $P_i$ and expanding in a neighborhood of $(v^s_i,v^u_i) = (0,0)$ gives the expansion~\eqref{Shil}. Moreover, this change of variable gives that $W^s_\mathrm{loc}(\tilde\gamma_i(0)) = \{v^u_i = 0\}$ and $W^u_\mathrm{loc}(\tilde\gamma_i(0)) = \{v^s_i = 0\}$. Finally, we may take $\delta := \min \{\delta_1,\dots,\delta_{k+1}\}$ to arrive at the $\delta > 0$ in the statement of the lemma. This completes the proof.
\end{proof} %End of proof

\end{document}
